\chapter{项目获取与初步判断}

\section{项目流的来源}

冷触达（cold outreach）在缺乏品牌时效率低但并非无效：关键是\textbf{高度定制}的切入点（对目标公司近期融资、产品或监管变化的具体评论）。更可持续的来源包括：创始人推荐、同行 FA、律师与会计师转介、行业会议与学术合作。\textbf{网络效应}随成功案例累积——第一单往往最难，需在冷启动章节所述触达策略与这里所述来源之间建立闭环。

\section{Teaser 与保密}

卖方 Teaser 通常为匿名一页纸，概述行业、财务区间与交易诉求。买方在签署\textbf{保密协议（NDA）}前应明确：接触信息的人员范围、用途限制、存续期与管辖法律。NDA 不是万能盾：敏感信息仍应分层披露，核心技术细节可推迟至 exclusivity 阶段。

\section{商业逻辑的快速体检}

在完整尽调之前，用有限时间回答：\textbf{客户为什么买单？}替代方案与切换成本？\textbf{单位经济模型}是否随规模改善？竞争格局是赢家通吃还是分散？若上述任一项无法在假设表格中自洽，可暂缓投入重度资源。市场大小（TAM/SAM/SOM）宜用多种方法交叉校验，避免单一夸张口径。

\section{与内部流程的衔接}

立项会应产生\textbf{书面备忘录}：投资主题、主要风险、下一步工作清单与负责人。避免「口头立项、无编号」导致后续 DD 与法务脱节。对于可能触发利益冲突的项目（如 LP 关联），应尽早启动合规标注与回避程序。
